% Part: turing-machines
% Chapter: machines-computations
% Section: configuration

\documentclass[../../../include/open-logic-section]{subfiles}

\begin{document}

\olfileid{tur}{mac}{con}
\olsection{پیکربندی‌ها و محاسبه‌ها}

\begin{explain}
دنبال‌کردن پیکربندی‌های ماشین زوج را در بالا به یاد آورید. سازوکار
خیالیِ متشکل از نوار، سر خواندن--نوشتن و برنامهٔ ماشین تورینگ در واقع
تنها راهی شهودی برای مجسم‌کردن محاسبهٔ ماشین تورینگ است. به‌طور رسمی
می‌توانیم محاسبهٔ ماشین تورینگ روی ورودی داده‌شده را دنباله‌ای از
\emph{پیکربندی‌ها} تعریف کنیم---و خود پیکربندی دنباله‌ای از نمادها
(متناظر با محتوای نوار در نقطه‌ای از محاسبه)، عددی نشان‌دهندهٔ جای سر
خواندن--نوشتن، و یک حالت است. با این‌ها می‌توانیم تعریف کنیم ماشین
تورینگ~$M$ روی ورودی داده‌شده چه چیزی را محاسبه می‌کند.
\end{explain}

\begin{defn}[پیکربندی]
\emph{پیکربندیِ} ماشین تورینگ
$M=\tuple{Q,\Sigma,q_0,\delta}$ سه‌تاییِ $\tuple{C,m,q}$ است که در آن
\begin{enumerate}
\item $C\in\Sigma^*$ دنباله‌ای متناهی از نمادهای $\Sigma$ است؛
\item $m\in\Nat$ عددی است که $<\len{C}$؛ و
\item $q\in Q$.
\end{enumerate}
به‌طور شهودی، دنبالهٔ~$C$ محتوای نوار است (نمادهای همهٔ خانه‌ها از
چپ‌ترین خانه تا آخرین خانهٔ ناتهی یا خانه‌ای که پیش‌تر دیده شده)،
$m$ شمارهٔ خانه‌ای است که سر خواندن--نوشتن می‌خواند (با شمارهٔ~$0$
برای چپ‌ترین خانه)، و $q$ حالت کنونی ماشین است.
\end{defn}

\begin{explain}
ورودی ممکن برای ماشین تورینگ دنباله‌ای از نمادهاست، معمولاً دنباله‌ای
که عددی را به شکلی کد می‌کند. پیکربندی آغازین ماشین تورینگ همان
پیکربندی‌ای است که ماشین را در آن برای کار روی آن ورودی آغاز می‌کنیم:
نوار نشانگر انتهای نوار را دارد و بلافاصله پس از آن ورودی روی خانه‌های
سمت راست نوشته شده است؛ سر خواندن--نوشتن چپ‌ترین خانهٔ ورودی (یعنی
خانهٔ سمت راست نشانگر انتهای چپ) را می‌خواند؛ و سازوکار در حالت آغازین
تعیین‌شدهٔ~$q_0$ است.
\end{explain}

\begin{defn}[پیکربندی آغازین]
\emph{پیکربندی آغازینِ} $M$ برای ورودی $I\in\Sigma^*$ برابر است با
\[
\tuple{\TMendtape \frown I \frown \TMblank, 1, q_0}.
\]
\end{defn}

\begin{explain}
نماد~$\frown$ برای \emph{الحاق} است---رشتهٔ ورودی بلافاصله پس از
نشانگر انتهای چپ آغاز می‌شود.
\end{explain}

\begin{defn}
می‌گوییم پیکربندی $\tuple{C,m,q}$ (برحسب~$M$)
\emph{در یک گام پیکربندی $\tuple{C',m',q'}$ را به دست می‌دهد} اگر و
تنها اگر
\begin{enumerate}
\item نماد $m$امِ $C$ برابر $\sigma$ باشد؛
\item مجموعهٔ دستورهای $M$ مشخص کند
  $\delta(q,\sigma)=\tuple{q',\sigma',D}$؛
\item نماد $m$امِ $C'$ برابر $\sigma'$ باشد؛ و
\item
\begin{enumerate}
\item $D=L$ و اگر $m>0$ آنگاه $m'=m-1$، وگرنه $m'=0$؛ یا
\item $D=R$ و $m'=m+1$؛ یا
\item $D=N$ و $m'=m$؛
\end{enumerate}
\item اگر $m'=\len{C}$، آنگاه $\len{C'}=\len{C}+1$ و نماد $m'$امِ
  $C'$ برابر~$\TMblank$ است. وگرنه $\len{C'}=\len{C}$؛
\item به ازای همهٔ $i$ها که $i<\len{C}$ و $i\neq m$،
  $C'(i)=C(i)$.
\end{enumerate}
\end{defn}

\begin{defn}\ollabel{defn:run-output}
\emph{اجرای $M$ روی ورودی~$I$} دنباله‌ای متناهی یا نامتناهی چون $C_i$
از پیکربندی‌های $M$ است، که $C_0$ پیکربندی آغازین $M$ برای ورودی~$I$
است و هر $C_i$ که جملهٔ پایانی دنباله نیست، در یک گام $C_{i+1}$ را
به دست می‌دهد.

می‌گوییم $M$ \emph{پس از $k$ گام روی ورودی $I$ متوقف می‌شود} اگر
$C_k=\tuple{C,m,q}$، نماد $m$امِ~$C$ برابر~$\sigma$ و
$\delta(q,\sigma)$ تعریف‌نشده باشد. در این حالت \emph{خروجیِ}~$M$
برای ورودی~$I$ برابر~$O$ است، که در آن $O$ رشته‌ای از نمادهاست که به
$\TMblank$ پایان نمی‌یابد و برای یک~$j\in\Nat$،
$C=\TMendtape\concat O\concat\TMblank^j$.
($\TMblank^j$ دنباله‌ای از $j$ عدد $\TMblank$ است.)
\end{defn}

\begin{explain}
بنا بر این تعریف، خروجی~$O$ از~$M$ همواره به نمادی غیر از
$\TMblank$ پایان می‌یابد؛ یا اگر در زمان~$k$ کل نوار جز چپ‌ترین
$\TMendtape$ با~$\TMblank$ پر باشد، $O$ رشتهٔ تهی است.
\end{explain}

\end{document}
